\documentclass[a4paper,11pt]{article}

\usepackage[a4paper,margin=2cm]{geometry}
\usepackage[T1]{fontenc}
\usepackage[utf8]{inputenc}
\usepackage[ngerman,english]{babel}
\usepackage{lmodern}
\usepackage{microtype}
\usepackage{xcolor}
\usepackage{parskip}
\usepackage{enumitem}
\usepackage[most]{tcolorbox}
\usepackage{titlesec}
\usepackage[hidelinks]{hyperref}

\definecolor{HeaderBlueDark}{RGB}{70,110,170}
\definecolor{RowBlueLight}{RGB}{235,243,255}
\definecolor{RowBlueDark}{RGB}{220,233,250}
\definecolor{SoftGray}{RGB}{90,90,90}

\setlength{\parindent}{0pt}
\setlist[itemize]{leftmargin=1.4em,itemsep=0.35em,topsep=0.35em}
\linespread{1.06}

\titleformat{\section}[block]
{\Large\bfseries\color{HeaderBlueDark}}
{}{0pt}{}
\titlespacing{\section}{0pt}{1.3em}{0.8em}

\titleformat{\subsection}[block]
{\large\bfseries\color{HeaderBlueDark}}
{}{0pt}{}
\titlespacing{\subsection}{0pt}{1.0em}{0.6em}

\newtcolorbox{exbox}{enhanced,breakable,colback=RowBlueLight,colframe=RowBlueDark,boxrule=0.4pt,arc=1.5mm,left=1.2mm,right=1.2mm,top=1mm,bottom=1mm,title={Beispiele \textnormal{(Examples)}},fonttitle=\bfseries,colbacktitle=RowBlueDark,coltitle=black}

\NewDocumentEnvironment{grammarentry}{mm+b}{%
    \begin{tcolorbox}[enhanced,breakable,colback=white,colframe=HeaderBlueDark,boxrule=0.6pt,arc=1.5mm,left=1.5mm,right=1.5mm,top=1mm,bottom=1mm,colbacktitle=HeaderBlueDark,coltitle=white,fonttitle=\bfseries,title={#1\hfill\normalfont\footnotesize #2}]
    #3
    \end{tcolorbox}
}{}

\newcommand{\GermanText}[1]{\textbf{Deutsch: }#1\par}
\newcommand{\EnglishText}[1]{{\small\color{SoftGray}\textbf{English: }#1\par}}
\newcommand{\ExampleItem}[2]{\item \textbf{#1}\\{\small\color{SoftGray}#2}}

\title{Das Adverb \glqq dahin\grqq}
\date{}

\begin{document}
    \maketitle
    \tableofcontents
    \newpage
    
    
    \section{Grundbedeutung}
        
        \begin{grammarentry}{dahin}{Adverb, Lokaladverb, direktional}
            \GermanText{\textbf{dahin} bedeutet: zu diesem Ort, in diese Richtung oder zu diesem Ziel. Es beantwortet normalerweise die Frage \glqq Wohin?\grqq.}
            \EnglishText{\textbf{dahin} means: to that place, in that direction, or toward that destination. It normally answers the question ``Where to?''}
            
            \GermanText{Aussprache: Betonung auf der zweiten Silbe: da-HIN.}
            \EnglishText{Pronunciation: stress on the second syllable: da-HIN.}
            
            \GermanText{Das Wort ist nicht steigerbar.}
            \EnglishText{The word cannot be compared.}
        \end{grammarentry}
    
    
    \section{Grammatische Funktion}
        
        \begin{grammarentry}{Adverbart}{Art und Funktion}
            \GermanText{\textbf{dahin} ist ein direktionales Lokaladverb. Es bezeichnet keine feste Position, sondern eine Richtung oder ein Ziel.}
            \EnglishText{\textbf{dahin} is a directional adverb of place. It does not describe a fixed position, but rather a direction or destination.}
            
            \GermanText{Fragewort: \textbf{Wohin?}}
            \EnglishText{Question word: \textbf{Where to?}}
        \end{grammarentry}
    
    
    \section{Unterschied zwischen \glqq dort\grqq, \glqq dorthin\grqq und \glqq dahin\grqq}
        
        \begin{grammarentry}{Vergleich}{Ort und Richtung}
            \GermanText{\textbf{dort} beantwortet die Frage \glqq Wo?\grqq und bezeichnet einen Ort ohne Bewegung.}
            \EnglishText{\textbf{dort} answers the question ``Where?'' and describes a place without movement.}
            
            \GermanText{\textbf{dahin} und \textbf{dorthin} beantworten die Frage \glqq Wohin?\grqq und bezeichnen eine Richtung oder ein Ziel.}
            \EnglishText{\textbf{dahin} and \textbf{dorthin} answer the question ``Where to?'' and describe a direction or destination.}
            
            \GermanText{\textbf{dorthin} ist oft etwas genauer oder stärker betont als \textbf{dahin}. In vielen Sätzen sind beide möglich.}
            \EnglishText{\textbf{dorthin} is often slightly more precise or more strongly emphasized than \textbf{dahin}. In many sentences, both are possible.}
        \end{grammarentry}
        
        \begin{exbox}
            \begin{itemize}
                \ExampleItem{Ich bin dort.}{I am there.}
                \ExampleItem{Ich gehe dahin.}{I am going there.}
                \ExampleItem{Ich fahre dorthin, wo meine Eltern wohnen.}{I am going to the place where my parents live.}
            \end{itemize}
        \end{exbox}
    
    
    \section{Räumliche Verwendung}
        
        \begin{grammarentry}{Bewegung zu einem Ort}{wörtliche Bedeutung}
            \GermanText{Bei der räumlichen Verwendung bezeichnet \textbf{dahin} ein Ziel, auf das sich eine Person oder eine Sache zubewegt.}
            \EnglishText{In its spatial use, \textbf{dahin} describes a destination toward which a person or thing moves.}
        \end{grammarentry}
        
        \begin{exbox}
            \begin{itemize}
                \ExampleItem{Ich stelle die Tasche dahin.}{I put the bag there.}
                \ExampleItem{Wir gehen dahin, wo es ruhiger ist.}{We are going to where it is quieter.}
                \ExampleItem{Leg das Buch bitte dahin.}{Please put the book there.}
            \end{itemize}
        \end{exbox}
    
    
    \section{Übertragene Verwendung}
        
        \begin{grammarentry}{Entwicklung zu einem Ergebnis}{übertragene Bedeutung}
            \GermanText{\textbf{dahin} kann auch eine Entwicklung zu einem Zustand, Ergebnis oder einer Schlussfolgerung ausdrücken.}
            \EnglishText{\textbf{dahin} can also express a development toward a state, result, or conclusion.}
            
            \GermanText{Eine häufige Struktur ist: \textbf{sich dahin entwickeln, dass ...}}
            \EnglishText{A common structure is: \textbf{sich dahin entwickeln, dass ...} (to develop to the point that ...)}
        \end{grammarentry}
        
        \begin{exbox}
            \begin{itemize}
                \ExampleItem{Die Situation hat sich dahin entwickelt, dass wir handeln müssen.}{The situation has developed to the point that we must act.}
                \ExampleItem{Die Diskussion ging dahin, dass eine neue Lösung notwendig ist.}{The discussion moved toward the conclusion that a new solution is necessary.}
                \ExampleItem{Seine Gedanken gehen dahin, das Projekt zu beenden.}{His thoughts are moving toward ending the project.}
            \end{itemize}
        \end{exbox}
    
    
    \section{Wichtige Verbindungen}
        
        \begin{grammarentry}{Typische Kombinationen}{Kollokationen}
            \GermanText{\textbf{dahin gehen} = zu diesem Ort gehen}
            \EnglishText{\textbf{dahin gehen} = to go there}
            
            \GermanText{\textbf{dahin fahren} = zu diesem Ort fahren}
            \EnglishText{\textbf{dahin fahren} = to travel or drive there}
            
            \GermanText{\textbf{etwas dahin stellen oder legen} = etwas an diesen Ort stellen oder legen}
            \EnglishText{\textbf{etwas dahin stellen oder legen} = to put something there}
            
            \GermanText{\textbf{sich dahin entwickeln, dass ...} = sich bis zu einem bestimmten Ergebnis entwickeln}
            \EnglishText{\textbf{sich dahin entwickeln, dass ...} = to develop to the point that ...}
            
            \GermanText{\textbf{dahin kommen, dass ...} = einen Zustand oder eine Schlussfolgerung erreichen}
            \EnglishText{\textbf{dahin kommen, dass ...} = to reach a state or conclusion where ...}
        \end{grammarentry}
    
    
    \section{Abgrenzung zu einem Pronomen}
        
        \begin{grammarentry}{Kein Ersatz für ein Nomen}{wichtiger Hinweis}
            \GermanText{\textbf{dahin} ist kein Pronomen und ersetzt normalerweise kein Nomen wie \glqq die Frage\grqq. Es bezeichnet vielmehr die Richtung oder das Ziel einer Veränderung.}
            \EnglishText{\textbf{dahin} is not a pronoun and normally does not replace a noun such as ``the question.'' Instead, it describes the direction or goal of a change.}
            
            \GermanText{In dem Satz \glqq Die Frage hat sich dahin verändert, dass ...\grqq bezeichnet \textbf{dahin} das Ergebnis oder die Richtung der Veränderung.}
            \EnglishText{In the sentence ``The question changed to the point that ...'', \textbf{dahin} describes the result or direction of the change.}
        \end{grammarentry}
    
    
    \section{Synonyme und nahe Ausdrücke}
        
        \begin{grammarentry}{Ähnliche Wörter}{je nach Kontext}
            \GermanText{Mögliche ähnliche Ausdrücke sind: \textbf{dorthin}, \textbf{in diese Richtung}, \textbf{zu diesem Punkt}, \textbf{zu diesem Ziel}.}
            \EnglishText{Possible similar expressions are: \textbf{dorthin}, \textbf{in this direction}, \textbf{to this point}, \textbf{toward this goal}.}
            
            \GermanText{Diese Ausdrücke sind nicht in jedem Satz vollständig austauschbar.}
            \EnglishText{These expressions are not completely interchangeable in every sentence.}
        \end{grammarentry}
    
    
    \section{Zusammenfassung}
        
        \begin{grammarentry}{Merksatz}{kurz und wichtig}
            \GermanText{\textbf{dort} = Wo? \quad \textbf{dahin / dorthin} = Wohin?}
            \EnglishText{\textbf{dort} = Where? \quad \textbf{dahin / dorthin} = Where to?}
            
            \GermanText{\textbf{dahin} bezeichnet entweder eine räumliche Richtung oder eine übertragene Entwicklung zu einem Ergebnis.}
            \EnglishText{\textbf{dahin} describes either a spatial direction or a figurative development toward a result.}
        \end{grammarentry}

\end{document}
